\documentclass[11pt]{article}
\usepackage[margin=1in]{geometry}
\usepackage{amsmath,amssymb,amsthm}
\usepackage{enumitem}
\usepackage{xcolor}
\usepackage{parskip}
\usepackage{booktabs}
\usepackage{hyperref}
\hypersetup{colorlinks=true, urlcolor=blue!45!black, linkcolor=blue!45!black}
\title{\vspace{-2em}\textbf{Differential Expression \& Multiple-Testing Correction}\\[0.3em]\large p-values, Benjamini--Hochberg FDR, and q-values\\[0.6em]\normalsize\textit{Computational Genomics 76553 --- recap of Lecture 5}}
\author{}\date{}

\newcommand{\note}[2]{\par\medskip\noindent\fbox{\parbox{0.97\linewidth}{\textbf{\textcolor{blue!45!black}{#1}}\\[0.2em]\textit{#2}}}\par\medskip}

\begin{document}
\maketitle
\tableofcontents
\bigskip

\begin{center}
\textit{``Throw enough nets into the sea and you will catch plastic bottles, not just fish.''}
\end{center}

\noindent From a single p-value to honest large-scale discovery: p-values, Bonferroni, Benjamini--Hochberg, and q-values. A recap of Lecture 5 --- p-values, multiple hypothesis testing, and FDR. Following the lecture's framing: from a single anomalous coverage region to testing the whole genome.

\section{The setting: from one region to many tests}

Lecture 5 starts where the previous lectures left off. We had modelled genome coverage with a Poisson distribution: the probability of coverage starting at position $i$ over length $k$ is Poisson in the requested depth and length. That Poisson model is our null hypothesis $H_0$ --- it describes a typical region, where coverage is random and independent of position.

The biological question is: are there special regions? Regions that ``attract'' or ``repel'' coverage, that carry more errors, that simply do not look like the null predicts? Given some observed coverage at some sequencing depth, we want to ask: is this anomalous or not, and how do we quantify it?

\note{Recap from CBIO}{There we discussed separating two simple hypotheses --- likelihood-ratio tests, Neyman--Pearson. Today is different: we have only $H_0$, and no alternative. We are not asking ``which of two hypotheses is more likely?'' but ``how unlikely is what we saw, assuming $H_0$?'' That single question is what the p-value answers.}

\subsection{What ``differential expression'' is, in this language}

In a single-cell / bulk genomics context, the very same machinery powers differential gene expression (DE). For each gene we ask whether its expression differs between two conditions or groups (e.g.\ treated vs.\ control, cell-type A vs.\ B). The per-gene null and alternative are:
\begin{align*}
H_{0,g} &: \text{gene } g \text{ is \emph{not} differentially expressed (same distribution in both groups)}\\
H_{1,g} &: \text{gene } g \text{ is differentially expressed}
\end{align*}

We compute a per-gene test statistic, turn it into a p-value, and --- crucially --- we do this for \emph{every} gene at once ($\sim$20{,}000 of them). That is exactly the ``many regions in the genome'' problem the lecture builds toward, and it is where naive p-value thresholds fall apart.

\section{The p-value}

Start with the abstract definition the lecture uses. We have a random variable $X$ and a null distribution $P_0$ on $X$. In our case $X$ is the coverage and $P_0$ is the Poisson model. We observe a single value $x$ and ask how surprising it is under $P_0$. The (right-tail) p-value is:
\[
\mathrm{pV}(X) \;=\; P_0\!\left(X \ge x\right) \;=\; P\!\left(\text{stat} \ge t \mid H_0\right).
\]

Intuitively: the chance of getting the observed value or anything more extreme. The lecture is careful here --- we use a tail probability, not a point probability. A single point can have tiny (or even zero, in the continuous case) probability while the whole region is perfectly ordinary. The tail is what measures ``extremeness.''

\note{When is a result surprising?}{When the p-value is small. We fix a significance threshold $\alpha$, and declare: if $\mathrm{pV} < \alpha$, the result is ``significant'' --- meaning significantly abnormal, i.e.\ $x$ is an outlier. How unlikely? ``$\alpha$-unlikely.''}

\subsection{Worked example from the lecture}

How surprising is coverage of $52$ on a depth of $D = 10$? We compute the tail probability under the Poisson null,
\[
\mathrm{pV} \;=\; P_0(X \ge 52) \;=\; 1 - \sum_{j=0}^{51} \frac{e^{-10}\,10^{\,j}}{j!}.
\]
The complement / CDF form is easier to compute, and the answer is: getting coverage $52$ when the expected depth is $10$ is extremely unlikely --- a tiny p-value, a clear outlier.

\note{What a p-value is \emph{not}}{A p-value is $P(\text{data this extreme} \mid H_0)$. It is \emph{not} $P(H_0 \text{ is true} \mid \text{data})$. The null is not a random hypothesis --- it is the fixed baseline (``no effect'', ``nothing special''). To go from $P(\text{data}\mid H_0)$ to $P(H_0 \mid \text{data})$ you would need a prior on $H_0$ and Bayes' rule. The p-value is strictly one-directional.}

\subsection*{Two things to always keep in mind}
\begin{itemize}[leftmargin=1.4em]
\item $\mathrm{pV}$ is a function of $H_0$. Always ask: which null am I testing against? Write it as $\mathrm{pV}(X; H_0)$.
\item $\mathrm{pV}$ always involves a tail (cumulative) probability. For Poisson or geometric this is trivial; for many distributions the tail is exactly the hard part, needing approximations.
\end{itemize}

\note{A subtle but central fact}{If the test statistic is continuous, then under $H_0$ the p-value is uniform on $[0,1]$:
\[
P_0\!\left(\mathrm{pV}(X) \le \rho\right) = \rho.
\]
Reason: $\mathrm{pV} \le \rho$ is exactly the event ``my statistic landed in the top $\rho$-fraction of the null tail'', and under the null the statistic is a draw from that distribution, so it lands there exactly a $\rho$-fraction of the time. This calibration is the foundation everything below rests on. (For discrete statistics you get $\le \rho$ instead of equality.)}

\section{The multiple-testing problem}

Now we change perspective: instead of one region with high coverage, we scan many regions across the genome and ask which look surprising. We get a whole vector of p-values, one per region:
\[
\mathrm{pV}_1,\ \mathrm{pV}_2,\ \dots,\ \mathrm{pV}_m
\qquad\longrightarrow\qquad
m \text{ separate null hypotheses.}
\]

What do we expect? Suppose we use $\alpha = 0.05$. The chance any single region is flagged anomalous by chance is $1$ in $20$. But now we test, say, $20{,}000$ hypotheses at once. The number flagged is a sum of indicators $\mathbf{1}\{\mathrm{pV}_i < \alpha\}$, and by linearity of expectation:
\[
\mathbb{E}[\#\text{ flagged}] \;=\; \sum_{i=1}^{m} P(\mathrm{pV}_i < \alpha) \;=\; m \cdot \alpha \;=\; 20{,}000 \cdot 0.05 \;=\; 1{,}000.
\]

\note{The core danger}{Even if nothing interesting is happening --- every null is true --- running $20{,}000$ tests at $\alpha = 0.05$ produces about $1{,}000$ false positives by pure chance. Roll a die enough times and even ``$1$-in-$20$'' events happen constantly. Real signal gets buried in noise. This is why so many published ``extraordinary'' findings are not extraordinary at all: with enough tests, something striking is guaranteed to appear (the ``look-elsewhere'' effect).}

The central question of the rest of the lecture: how do we correct the framework so we don't get $\alpha \cdot m$ spurious ``discoveries'' when the null is true? Since the false-positive count is linear in $\alpha$ and in the number of tests, every correction revolves around shrinking or reinterpreting $\alpha$.

\begin{center}
\begin{tabular}{@{}lll@{}}
\toprule
\textbf{Quantity to control} & \textbf{Name} & \textbf{Plain meaning}\\
\midrule
$P(\text{at least one false positive})$ & FWER (family-wise error rate) & ``I refuse to report even one wrong result.''\\[0.3em]
$\mathbb{E}[\,\text{false pos.} / \text{total called}\,]$ & FDR (false discovery rate) & ``A small fraction of my reported hits may be wrong.''\\
\bottomrule
\end{tabular}
\end{center}

\section{Bonferroni: control the family-wise error}

The Bonferroni correction takes a corrected threshold:
\[
\alpha_{\mathrm{Bonf}} \;=\; \frac{\alpha}{m}.
\]

Why does this work at scale? We want to bound the probability that even one of the $m$ tests drops below the threshold under $H_0$. This is a question about a union of events, so we apply the union bound:
\[
P\!\left(\exists\, i : \mathrm{pV}_i \le \tfrac{\alpha}{m}\right)
\;\le\; \sum_{i=1}^{m} P\!\left(\mathrm{pV}_i \le \tfrac{\alpha}{m}\right)
\;=\; m \cdot \frac{\alpha}{m}
\;=\; \alpha.
\]

\note{What Bonferroni buys}{The correction restores $\alpha$ for the whole system: the probability of getting even one spurious hit at level $0.05$ across all $m$ tests is back to $0.05$. If you want a family-wise error of $\alpha$ over $m$ tests, test each at $\alpha/m$. The more questions you ask, the smaller the per-test threshold.}

\note{The problem with Bonferroni}{It asks too strict a question. It controls the chance that \emph{anyone at all} crosses the threshold --- it tries to guarantee almost no false positive ever. In real genomics, with $20{,}000$ genes, $\alpha/m$ becomes so tiny (e.g.\ $0.05/20000 = 0.0000025$) that you lose almost all real biological signal. We usually do not need ``no false positives ever''; we need to control the \emph{expected fraction} of false hits among those we report.}

\section{FDR and the Benjamini--Hochberg procedure}

Why is FDR a better fit for multiple testing? Because the population of tests usually hides two groups: a large majority that genuinely follow $H_0$, and a small hidden group of truly significant effects. Bonferroni, in its zeal, misses the small real group. FDR instead controls the \emph{rate} of false discoveries rather than forbidding them.

\note{The fishing-net intuition}{Throw a net into the sea. You may catch plastic bottles instead of fish. We don't demand that every cast contains only fish (impossible). We demand that \emph{most} of what we haul up is fish. FDR controls the share of junk in the catch, not the existence of junk.}

\subsection{Defining FDR}

Among the regions/genes we call positive, the FDR is the expected fraction that are actually false positives:
\[
\mathrm{FDR} \;=\; \mathbb{E}\!\left[\frac{V}{\max(R,\,1)}\right],
\qquad V = \#\text{false positives}, \quad R = \#\text{total called}.
\]

Controlling FDR at $0.05$ says: of the things I report, on average at most $5\%$ are noise. Unlike Bonferroni (``$1$-in-$20$ chance of \emph{anything} spurious''), FDR says ``take the whole pile of hits, and at most $1$-in-$20$ of \emph{them} is wrong.''

\subsection{Deriving the BH threshold}

Sort the p-values from most to least surprising and consider declaring the first $k$ of them discoveries:
\[
Y_{(1)} \le Y_{(2)} \le \dots \le Y_{(m)}
\qquad\longrightarrow\qquad
\text{cut after } k, \text{ using threshold } p = Y_{(k)}.
\]

Under the null, $P(\text{p-value} \le p) = p$, so among $m$ tests the expected number of false discoveries at threshold $p$ is at most $m \cdot p$. With $k$ discoveries total, the expected false fraction is bounded by $m \cdot p / k$. Demanding this be $\le \alpha$ and substituting $p = Y_{(k)}$:
\[
\frac{m \cdot Y_{(k)}}{k} \;\le\; \alpha
\qquad\Longrightarrow\qquad
Y_{(k)} \;\le\; \frac{k}{m}\,\alpha.
\]

\note{Benjamini--Hochberg procedure (control FDR at level $q$)}{%
\begin{enumerate}[leftmargin=1.6em,nosep]
\item Sort all $m$ p-values in increasing order: $\;Y_{(1)} \le Y_{(2)} \le \dots \le Y_{(m)}$.
\item For each rank $k$, compare to the BH line: $\;Y_{(k)} \le \dfrac{k}{m}\,q\;$?
\item Find the \textbf{largest} $k$ for which this holds. \quad ($\leftarrow$ the cutoff)
\item Declare \emph{all} tests with rank $\le k$ significant (even any whose own p-value is above its own line).
\end{enumerate}}

\note{Why the moving threshold?}{Bonferroni uses one tiny fixed threshold $\alpha/m$ for every test. BH uses $\dfrac{k}{m}\,\alpha$, which \emph{grows} with $k$. At $k=1$ it equals Bonferroni; for larger $k$ it relaxes. The logic: if you are already making many discoveries, tolerating a few false ones is fine as long as their fraction stays below $\alpha$.}

\section{Worked BH mini-example}

Take $m = 5$ tests with ordered p-values, and control FDR at $\alpha = 0.05$. The BH line at rank $k$ is $\dfrac{k}{5}\cdot 0.05$.

\begin{center}
\begin{tabular}{@{}ccccc@{}}
\toprule
Rank $k$ & p-value $Y_{(k)}$ & BH line $\frac{k}{5}\cdot 0.05$ & $Y_{(k)} \le$ line? & ``needed FDR'' $= \frac{m\,Y_{(k)}}{k}$\\
\midrule
$1$ & $0.002$ & $0.010$ & yes & $0.010$\\
$2$ & $0.010$ & $0.020$ & yes & $0.025$\\
$3$ & $0.030$ & $0.030$ & yes $\leftarrow$ largest $k$ & $0.050$\\
$4$ & $0.200$ & $0.040$ & no & $0.250$\\
$5$ & $0.600$ & $0.050$ & no & $0.600$\\
\bottomrule
\end{tabular}
\end{center}

The largest $k$ satisfying $Y_{(k)} \le \dfrac{k}{m}\,\alpha$ is $k = 3$. So we declare ranks $1, 2, 3$ as discoveries. Note this is a step-up rule: we find the rightmost crossing and accept everything to its left, even if an intermediate point sat just on or below its own line.

\note{Reading the last column}{The final column, $\dfrac{m\,Y_{(k)}}{k}$, is essentially the \emph{q-value} of each test: the smallest FDR level at which that test would be called significant. Result 1 needs only FDR $1\%$; result 3 needs FDR $5\%$; result 4 would need FDR $25\%$. (A small monotonicity fix makes q-values formally non-decreasing, but this is the idea.)}

\section{q-values and how significance is reported}

Raw p-values are calibrated for one hypothesis. But in large-scale testing the question we actually care about is different: if I put this result on my discovery list, how reliable is that list? A p-value of $0.01$ across $20{,}000$ genes is not impressive --- under pure noise you expect $20{,}000 \cdot 0.01 = 200$ such hits. So we report a \emph{q-value} instead: a p-value recalibrated for the whole search.

\note{q-value, defined}{The q-value of a test is the smallest FDR level at which that test would be called significant by the BH procedure. Equivalently:
\[
q_{(k)} \;=\; \min_{k' \ge k}\ \frac{m\,Y_{(k')}}{k'},
\]
among all discoveries at least this significant, the expected fraction that are false. The p-value is to a single test what the q-value is to FDR.}

\begin{center}
\begin{tabular}{@{}lll@{}}
\toprule
 & \textbf{p-value} & \textbf{q-value}\\
\midrule
Question answered & How extreme is this under $H_0$? & What FDR must I tolerate to keep this on the list?\\
Scope & One hypothesis & Whole family of tests\\
Reporting cutoff & $p < 0.05$ (single test) & $q < 0.05$ (standard in genomics papers)\\
\bottomrule
\end{tabular}
\end{center}

BH and q-values solve the same problem from opposite ends: BH fixes $\alpha$ first and returns the discovery set; the q-value attaches to each test the minimal $\alpha$ that would include it. Reporting q-values lets a reader pick their own tolerance --- accept $q < 0.05$ for a strict list, $q < 0.10$ for a broader one --- without rerunning the procedure.

\subsection{Refinements the lecture flags}

The bound ``expected false discoveries $\le m \cdot p$'' pessimistically assumes \emph{all} hypotheses are null. If only $m_0$ are truly null, the real expectation is $m_0 \cdot p < m \cdot p$. Adaptive methods estimate the null proportion $\pi_0 = m_0/m$ and use $\pi_0 \cdot m \cdot p$ instead, recovering more power (Storey's q-value). The plain BH bound is tightest when $m$ is large and most tests are null.

\note{Effect size matters alongside p/q}{A tiny q-value says a difference is \emph{real}, not that it is \emph{large}. In DE analysis, always report the effect size (e.g.\ log fold-change) next to the p/q-value. A gene can be highly significant yet biologically trivial; volcano plots exist precisely to show both axes at once.}

\note{Assumptions \& test choice (DE / single-cell)}{BH controls FDR under independence or positive dependence of the tests. Heavily / negatively correlated tests need Benjamini--Yekutieli (more conservative). The p-values themselves must be calibrated: a wrong null (e.g.\ Poisson when counts are over-dispersed) distorts interpretation --- negative-binomial models or empirical nulls fit count data better. Choose the per-gene test to match the data; FDR correction is only as honest as the p-values feeding it.}

\note{Core ideas to remember}{%
\begin{enumerate}[leftmargin=1.6em,nosep]
\item p-value $= P(\text{data this extreme} \mid H_0)$, not $P(H_0 \text{ true} \mid \text{data})$.
\item $\sim$20{,}000 tests at $\alpha = 0.05 \to \sim$1{,}000 false positives by chance --- correction is mandatory.
\item Bonferroni ($\alpha/m$) controls FWER --- safe but loses real signal at scale.
\item BH controls FDR via the largest $k$ with $Y_{(k)} \le \dfrac{k}{m}\,\alpha$.
\item q-value $=$ smallest FDR at which a test is called; report $q < 0.05$ with an effect size.
\end{enumerate}}

\end{document}
